% ============ 02. Contexto y guía de lectura ============
\section{Contexto y guía de lectura}

\subsection{Qué es este documento}

Este documento desarrolla el \textbf{criterio 6 de la rúbrica}:
\emph{Línea de tiempo de commits y construcción del repositorio}. El repositorio auditado es el
Observatorio de investigadores reconocidos del Ministerio de Ciencia,
Tecnología e Innovación de Colombia (MinCiencias), que consolida y analiza
\textbf{6 convocatorias (2013--2021)} del padrón oficial de investigadores
y su producción científica. Esta serie cubre los \textbf{6 criterios} de la
rúbrica en documentos separados, más el documento maestro que los integra.

\textbf{En resumen:} Este criterio documenta cómo se construyó el repositorio a lo largo del tiempo: 133 commits entre agosto de 2025 y mayo de 2026 por 11 autores, con la autoría de cada uno. Se distinguen tres momentos: fase académica inicial (ago-nov 2025), un hiato de \textasciitilde{}100 días sin commits (nov 2025 - feb 2026) y la fase de ingeniería y cierre (feb-may 2026).




\begin{tcolorbox}[nota]
\textbf{Cómo leer este documento:} cada PDF de esta serie desarrolla \emph{un}
criterio de la rúbrica. Los demás criterios se tratan a fondo en sus propios
documentos y en el \textbf{documento maestro} (\texttt{maestro\_auditoria.pdf}).
Si este es tu primer acercamiento, lee primero la portada y este contexto; al
final encontrarás las conclusiones del criterio.
\end{tcolorbox}

\subsection{Glosario (solo los términos que usa este documento)}

Esta tabla incluye únicamente los términos técnicos que aparecen en este
documento, explicados en lenguaje sencillo:

\begingroup
\small
\setlength{\tabcolsep}{3pt}
\begin{longtable}{>{\raggedright\arraybackslash}p{4.3cm}>{\raggedright\arraybackslash}p{9.8cm}}
\caption{Glosario de términos presentes en este documento}
\label{tab:glosario}\\
\toprule
\textbf{Término} & \textbf{Qué significa} \\
\midrule
\endfirsthead
\toprule
\textbf{Término} & \textbf{Qué significa} \\
\midrule
\endhead
    HHI (Herfindahl-Hirschman) & Índice que mide si algo está concentrado o repartido. Va de 0 (todo repartido) a 1 (todo en un solo lugar). En el informe se usa para ver si los investigadores se concentran en pocos departamentos.
 \\
    Matriz de transición & Tabla que muestra con qué probabilidad un investigador pasa de una categoría (p. ej. Junior) a otra (p. ej. Asociado) entre dos convocatorias. Permite ver si el sistema 'sube', 'mantiene' o 'pierde' investigadores.
 \\
    Modelo dimensional (esquema estrella) & Forma de organizar los datos en tablas de 'hechos' (métricas: quién, cuándo, qué categoría) y 'dimensiones' (descripciones: área, región, institución), para hacer consultas rápidas y reportes.
 \\
    DuckDB & Base de datos analítica ligera que se ejecuta dentro del propio proyecto (sin instalar un servidor). Se usa para almacenar el modelo dimensional.
 \\
    Pipeline & Cadena de procesos que lleva los datos desde la fuente hasta el análisis (ingesta, limpieza, transformación, modelado). También llamada 'cadena de procesamiento' o 'flujo de datos'.
 \\
    Sprint & Ciclo corto de trabajo (en este proyecto, de 2 semanas) con objetivos concretos.
 \\
\bottomrule
\end{longtable}
\endgroup
